% preambulo.tex — Preámbulo LaTeX para la tesis doctoral (Quarto → PDF)
% Configuración APA: tablas/figuras con título arriba y fuente abajo,
% citas con márgenes simétricos, IBM Plex como familia tipográfica.

% ─── Fuentes IBM Plex ─────────────────────────────────────────────
% La configuración de mainfont/monofont/sansfont va en _quarto.yml.
% Aquí configuramos aspectos complementarios.
\usepackage{fontspec}

% ─── Márgenes (APA doctoral) ──────────────────────────────────────
\usepackage[
  a4paper,
  top=2.5cm,
  bottom=2.5cm,
  left=3cm,
  right=2.5cm
]{geometry}

% ─── Interlineado ─────────────────────────────────────────────────
\usepackage{setspace}
\setstretch{1.5}

% ─── Block quotes: márgenes simétricos de 1cm extra ──────────────
\usepackage{csquotes}
\renewenvironment{quote}%
  {\list{}{\leftmargin=1cm\rightmargin=1cm\itshape}%
   \item\relax\small}%
  {\endlist}

% ─── Tablas: título arriba (APA) ─────────────────────────────────
\usepackage{caption}
\captionsetup[table]{
  position=above,
  labelfont={bf,small},
  textfont={small},
  labelsep=period,
  justification=raggedright,
  singlelinecheck=false,
  skip=6pt
}
\captionsetup[figure]{
  position=above,
  labelfont={bf,small},
  textfont={small},
  labelsep=period,
  justification=raggedright,
  singlelinecheck=false,
  skip=6pt
}

% Nota de fuente debajo de tablas (uso manual: \tablenote{...})
\newcommand{\tablenote}[1]{%
  \par\vspace{4pt}%
  {\small\textit{Nota.} #1}%
}

% ─── Encabezados (jerarquía visual APA-compatible) ────────────────
\usepackage{titlesec}

% Nivel 1: Capítulo — IBM Plex Sans Bold, grande
\titleformat{\chapter}[display]
  {\normalfont\huge\bfseries\sffamily}
  {\chaptertitlename\ \thechapter}{20pt}{\Huge}
\titlespacing*{\chapter}{0pt}{-20pt}{40pt}

% Nivel 2: Sección
\titleformat{\section}
  {\normalfont\Large\bfseries\sffamily}
  {\thesection}{1em}{}

% Nivel 3: Subsección
\titleformat{\subsection}
  {\normalfont\large\bfseries\sffamily}
  {\thesubsection}{1em}{}

% Nivel 4: Subsubsección (itálica bold, APA nivel 3)
\titleformat{\subsubsection}
  {\normalfont\normalsize\bfseries\itshape\sffamily}
  {\thesubsubsection}{1em}{}

% Nivel 5: Párrafo (bold, APA nivel 4)
\titleformat{\paragraph}
  {\normalfont\normalsize\bfseries\sffamily}
  {\theparagraph}{1em}{}

% Nivel 6: Subpárrafo (itálica)
\titleformat{\subparagraph}
  {\normalfont\normalsize\itshape\sffamily}
  {\thesubparagraph}{1em}{}

% ─── Hipervínculos ────────────────────────────────────────────────
\usepackage{hyperref}
\hypersetup{
  colorlinks=true,
  linkcolor={blue!60!black},
  citecolor={blue!60!black},
  urlcolor={blue!70!black},
  bookmarksnumbered=true,
  bookmarksopen=true
}

% ─── Líneas finas para tablas (estilo limpio) ─────────────────────
\usepackage{booktabs}

% ─── Colores personalizados ──────────────────────────────────────
\usepackage{xcolor}
\definecolor{accentblue}{HTML}{0F62FE}   % IBM Blue
\definecolor{darkgray}{HTML}{393939}

% ─── Notas al pie ────────────────────────────────────────────────
\usepackage[bottom,hang]{footmisc}
\setlength{\footnotemargin}{0.8em}

% ─── Código (snippets) con fondo gris ────────────────────────────
\usepackage{framed}
\definecolor{codebg}{HTML}{F4F4F4}
\AtBeginDocument{%
  \ifdefined\Shaded
    \renewenvironment{Shaded}{\begin{snugshade}\small}{\end{snugshade}}%
    \colorlet{shadecolor}{codebg}%
  \fi
}

% ─── Epígrafe / dedicatoria ──────────────────────────────────────
\newcommand{\epigraph}[2]{%
  \begin{flushright}
  \begin{minipage}{0.6\textwidth}
  \small\itshape #1
  \par\vspace{4pt}
  \upshape\small --- #2
  \end{minipage}
  \end{flushright}
  \vspace{1em}
}

% ─── Portadilla con directora (PDF) ──────────────────────────────
% Redefine \maketitle para incluir a la directora de tesis
% debajo del autor, en la página de título generada por Quarto.
\usepackage{etoolbox}
\makeatletter
\patchcmd{\maketitle}
  {\@author}
  {\@author\\[0.8em]
   {\normalsize\normalfont
    Directora de tesis:\\[0.2em]
    \textbf{Dra.~María Gabriela Palazzo}}}
  {}{}
\makeatother
